%============================================================
% CH2 · P1 · PROBLEM 2 — ROUND TRIP II (CSES 1678)
%============================================================
\subsection{Problem 2.1.2: Round Trip II}
\problemlabel{2.1.2}{ssec:ch2-round-trip-ii}

\problemcard{2.1.2}{Round Trip II}{CSES 1678}{https://cses.fi/problemset/task/1678}{\exmedium\ (the directed twin of Problem~2.1.1 --- one word changes the algorithm)}{Directed cycle detection with three colours: only an edge to a \emph{gray} vertex closes a cycle; print it by reversing the parent walk}

\subsubsection*{1. Problem statement}

Byteland has $n$ cities and $m$ \textbf{one-way} flights. Design a round trip that starts in a city, visits one or more other cities and returns to the start, with all intermediate cities distinct. Print the number $k$ of cities on the route (start counted twice) and the route, or \texttt{IMPOSSIBLE}. Any valid route is accepted.

\textbf{Constraints:} $1 \le n \le 10^5$, $1 \le m \le 2\cdot10^5$; every flight goes from a city to a \emph{different} city. Two flights $a \to b$ and $b \to a$ may both exist.

\begin{table}[H]
\centering\small
\begin{tabular}{@{}lll@{}}
\toprule
\textbf{Input} & \textbf{Official output} & \textbf{Our output} \\
\midrule
\texttt{4 5 / 1 3 / 2 1 / 2 4 / 3 2 / 3 4} & \texttt{4 / 2 1 3 2} & \texttt{4 / 1 3 2 1} \\
\bottomrule
\end{tabular}
\end{table}

\begin{diagrambox}
\begin{tikzpicture}[x=1.2cm, y=1cm]
    \begin{scope}
        \node[gv] (a1) at (0,0) {1}; \node[gv] (a3) at (1.6,0.8) {3};
        \node[gv] (a2) at (1.6,-0.8) {2}; \node[gv] (a4) at (3.2,0) {4};
        \draw[gde] (a1) -- (a3); \draw[gde] (a2) -- (a1); \draw[gde] (a2) -- (a4);
        \draw[gde] (a3) -- (a2); \draw[gde] (a3) -- (a4);
        \node[font=\scriptsize\bfseries] at (1.6,-1.65) {the sample};
    \end{scope}
    \begin{scope}[xshift=5cm]
        \node[gvsel] (b1) at (0,0) {1}; \node[gvsel] (b3) at (1.6,0.8) {3};
        \node[gvsel] (b2) at (1.6,-0.8) {2}; \node[gv] (b4) at (3.2,0) {4};
        \draw[gpath, -{Stealth[length=3mm]}] (b1) -- (b3); \draw[gpath, -{Stealth[length=3mm]}] (b3) -- (b2);
        \draw[gpath, -{Stealth[length=3mm]}] (b2) -- (b1);
        \draw[gde] (b2) -- (b4); \draw[gde] (b3) -- (b4);
        \node[font=\scriptsize\bfseries] at (1.6,-1.65) {the only cycle: $1 \to 3 \to 2 \to 1$};
    \end{scope}
    \begin{scope}[xshift=10cm]
        \node[gv] (c1) at (0,0) {1}; \node[gv] (c3) at (1.6,0.8) {3};
        \node[gv] (c2) at (1.6,-0.8) {2}; \node[gvdst] (c4) at (3.2,0) {4};
        \draw[gde] (c3) -- (c4); \draw[gde] (c2) -- (c4);
        \node[note, text width=3.2cm] at (1.6,-1.55) {city 4 has two incoming flights but lies on no cycle};
    \end{scope}
\end{tikzpicture}
\figcap{The sample. Undirected, $\{2,3,4\}$ would be a triangle too --- but $2 \to 4 \leftarrow 3$ cannot be walked around.}
\end{diagrambox}

\subsubsection*{2. The graph model}

\begin{modelbox}[Graph Model --- Round Trip II]
\renewcommand{\arraystretch}{1.35}
\begin{tabularx}{\linewidth}{@{}>{\bfseries}l X@{}}
Vertices & the $n$ cities \\
Edges & directed, one per flight; no self-loops; $a \to b$ and $b \to a$ may coexist (a 2-cycle is a valid trip: ``one or more other cities'') \\
Weights & none \\
Question & find any directed cycle \\
Size & $|V| \le 10^5$, $|E| \le 2\cdot10^5$: \VE{}, iterative DFS \\
\end{tabularx}
\end{modelbox}

\subsubsection*{3. Key observations}

\begin{itemize}[leftmargin=*]
    \item \textbf{Visited is not enough.} City 4 is reached twice ($3 \to 4$ and $2 \to 4$) without any cycle. In a directed graph a revisit means ``two paths lead here'', not ``I came back''.
    \item \textbf{Gray is exactly right.} An edge to a vertex that is still on the current DFS path closes a cycle; by the cycle theorem (Lesson~\ref{ssec:ch2-cycles}) DFS meets such an edge iff any cycle exists.
    \item \textbf{No parent skip.} Unlike Problem~\ref{ssec:ch2-round-trip}, the edge $v \to u$ back to the parent is a different flight from $u \to v$ and genuinely forms the trip $u, v, u$.
    \item \textbf{Direction of printing.} Parent pointers point \emph{against} the flights; the walk from $u$ up to $v$ must be reversed.
\end{itemize}

\subsubsection*{4. Approaches}

\begin{bruteforcebox}
\textbf{Idea: test every city as a possible start.} City $s$ lies on a cycle iff $s$ can be reached again from one of its successors. Run a BFS from the successors of $s$; if it arrives at $s$, the BFS parents spell the trip.

\begin{lstlisting}[style=cpp]
#include <bits/stdc++.h>
using namespace std;

// Brute force: city s lies on a cycle iff s can be reached again from one of its successors.
// For every s run a BFS that starts at the successors of s; if it reaches s, rebuild the path.
int main() {
    int n, m;
    cin >> n >> m;
    vector<vector<int>> adj(n + 1);
    for (int i = 0; i < m; i++) {
        int a, b;
        cin >> a >> b;
        adj[a].push_back(b);
    }
    for (int s = 1; s <= n; s++) {
        vector<int> par(n + 1, 0);
        queue<int> q;
        for (int v : adj[s])
            if (!par[v]) { par[v] = s; q.push(v); }
        while (!q.empty() && !par[s]) {
            int u = q.front(); q.pop();
            for (int v : adj[u])
                if (!par[v]) { par[v] = u; q.push(v); }
        }
        if (!par[s]) continue;                        // s is on no cycle
        vector<int> cyc;                              // walk back from s along par
        int x = s;
        do { cyc.push_back(x); x = par[x]; } while (x != s);
        cyc.push_back(s);
        reverse(cyc.begin(), cyc.end());
        cout << cyc.size() << "\n";
        for (int y : cyc) cout << y << " ";
        cout << "\n";
        return 0;
    }
    cout << "IMPOSSIBLE\n";
}
\end{lstlisting}

As a bonus the BFS finds a \emph{shortest} cycle through $s$. On the sample, $s = 1$ succeeds at once and the program prints \texttt{4 / 1 3 2 1}.
\begin{complexitybox}
\textbf{Time:} \bigO{n(n+m)} --- on a DAG every one of the $10^5$ searches fails after up to $3\cdot10^5$ steps. \textbf{Space:} \bigO{n+m}. \textbf{Verdict:} \tle; a reference for testing.
\end{complexitybox}
\end{bruteforcebox}

\begin{algorithmbox}[Approach 2: DFS with three colours]
\textbf{Idea.} Colour 0 = never entered, 1 = on the current path, 2 = finished. An edge $u \to v$ with colour 1 at $v$ is a back edge: the trip is $v \to \dots \to u \to v$ along the parent pointers. Iterative, so a chain of $10^5$ flights is no problem.

\begin{lstlisting}[style=cpp]
#include <bits/stdc++.h>
using namespace std;

// Round Trip II (CSES 1678): print a directed cycle or IMPOSSIBLE.
// Iterative DFS with colours 0 white / 1 gray / 2 black; a gray target closes a cycle.
int main() {
    ios::sync_with_stdio(false);
    cin.tie(nullptr);
    int n, m;
    cin >> n >> m;
    vector<vector<int>> adj(n + 1);
    for (int i = 0; i < m; i++) {
        int a, b;
        cin >> a >> b;
        adj[a].push_back(b);
    }
    vector<int> color(n + 1, 0), par(n + 1, 0);
    vector<size_t> it(n + 1, 0);
    for (int s = 1; s <= n; s++) {
        if (color[s]) continue;
        color[s] = 1;
        vector<int> st(1, s);
        while (!st.empty()) {
            int u = st.back();
            if (it[u] == adj[u].size()) {              // all edges of u examined
                color[u] = 2;
                st.pop_back();
                continue;
            }
            int v = adj[u][it[u]++];
            if (color[v] == 0) {                       // tree edge
                color[v] = 1; par[v] = u;
                st.push_back(v);
            } else if (color[v] == 1) {                // back edge u -> v: cycle v -> ... -> u -> v
                vector<int> cyc;
                for (int x = u; x != v; x = par[x]) cyc.push_back(x);
                cyc.push_back(v);
                reverse(cyc.begin(), cyc.end());
                cyc.push_back(v);
                cout << cyc.size() << "\n";
                for (int x : cyc) cout << x << " ";
                cout << "\n";
                return 0;
            }                                          // black: forward / cross edge, ignore
        }
    }
    cout << "IMPOSSIBLE\n";
}
\end{lstlisting}

\begin{complexitybox}
\textbf{Time:} \VE. \textbf{Space:} \bigO{n + m}. \textbf{Verdict:} \ac{}. On a DAG with $10^5$ cities and $2\cdot10^5$ flights (a long chain plus random forward flights) it prints \texttt{IMPOSSIBLE} in a fraction of a second; on the chain $1 \to 2 \to \dots \to 10^5 \to 1$ it prints the full trip of $100\,001$ cities.
\end{complexitybox}
\end{algorithmbox}

\begin{invariantbox}[Why gray means cycle and black means nothing]
\textbf{Invariant.} The gray vertices are exactly the stack $x_0, \dots, x_k = u$, with an edge $x_i \to x_{i+1}$ for every $i$; a black vertex has all its out-edges examined and \emph{no cycle was found through them}.

\emph{Gray target.} $v = x_j$: the flights $x_j \to x_{j+1} \to \dots \to x_k \to x_j$ form a trip with distinct intermediate cities.

\emph{Black target.} Every vertex reachable from a black vertex $w$ is black too: when $w$ finished, all its successors were black or had just been finished below it; if some vertex reachable from $w$ were gray, the path to it would have produced a back edge before $w$ finished. So from $v$ there is no way back onto the current stack, and $u \to v$ cannot be part of a cycle through gray vertices. Ignoring it is safe.

\emph{Completeness.} If a cycle exists, the cycle theorem guarantees a back edge, so the search cannot end with \texttt{IMPOSSIBLE}. \hfill$\blacksquare$
\end{invariantbox}

\begin{alternativebox}[Approach 3: peel the sources (Kahn), then walk backwards]
Repeatedly delete cities with no incoming flight (Kahn's algorithm, Chapter~5). If every city gets deleted, the graph is acyclic. Otherwise, every survivor still has an incoming flight from another survivor --- it was never deleted because its in-degree never dropped to 0. Start anywhere among the survivors and keep stepping to \emph{some surviving predecessor}. With finitely many cities a city must repeat, and the repeated stretch, read in flight direction, is a cycle.

\begin{lstlisting}[style=cpp]
#include <bits/stdc++.h>
using namespace std;

// Alternative: peel vertices of in-degree 0 (Kahn). If everything is peeled, there is no cycle.
// Otherwise every survivor has an incoming edge from another survivor: walk BACKWARDS along such
// edges until a vertex repeats; the repeated part, read forwards, is a cycle.
int main() {
    ios::sync_with_stdio(false);
    cin.tie(nullptr);
    int n, m;
    cin >> n >> m;
    vector<vector<int>> adj(n + 1), radj(n + 1);
    vector<int> indeg(n + 1, 0);
    for (int i = 0; i < m; i++) {
        int a, b;
        cin >> a >> b;
        adj[a].push_back(b);
        radj[b].push_back(a);
        indeg[b]++;
    }
    vector<char> alive(n + 1, 1);
    queue<int> q;
    for (int v = 1; v <= n; v++)
        if (indeg[v] == 0) q.push(v);
    int peeled = 0;
    while (!q.empty()) {
        int v = q.front(); q.pop();
        alive[v] = 0; peeled++;
        for (int w : adj[v])
            if (--indeg[w] == 0) q.push(w);
    }
    if (peeled == n) { cout << "IMPOSSIBLE\n"; return 0; }
    int start = 1;
    while (!alive[start]) start++;
    vector<int> pos(n + 1, -1), walk;                 // pos[v] = index of v in walk
    int v = start;
    while (pos[v] == -1) {
        pos[v] = walk.size();
        walk.push_back(v);
        for (int u : radj[v])                          // any alive predecessor will do
            if (alive[u]) { v = u; break; }
    }
    vector<int> cyc(walk.begin() + pos[v], walk.end());   // v <- ... <- (back to v)
    reverse(cyc.begin(), cyc.end());                       // now in edge direction
    cyc.insert(cyc.begin(), v);                            // v -> ... -> v
    cout << cyc.size() << "\n";
    for (int x : cyc) cout << x << " ";
    cout << "\n";
}
\end{lstlisting}
On the sample no city has in-degree 0, so nothing is peeled; the backward walk is $1 \leftarrow 2 \leftarrow 3 \leftarrow 1$ and the output is again \texttt{4 / 1 3 2 1}.
\begin{complexitybox}
\textbf{Time:} \VE{} (the backward walk scans each city's predecessor list at most once). \textbf{Space:} \bigO{n+m}. \textbf{Verdict:} \ac. No DFS at all, no recursion; the peeled order is a topological order of the deleted part, which Chapter~5 develops.
\end{complexitybox}
\end{alternativebox}

\begin{tricksbox}[Approach 4: Implementation Refinements]
\begin{itemize}[leftmargin=*]
    \item \textbf{One colour array for all starts.} Black vertices from an earlier start are never explored again; a fresh array per start would make the search quadratic.
    \item \textbf{Stop at the first cycle.} Returning immediately keeps the code short; the colours of unfinished vertices are never needed again.
    \item \textbf{Recursive version.} Possible, but a chain of $10^5$ flights makes it $10^5$ deep (Lesson~\ref{ssec:ch2-call-stack}).
    \item \textbf{Which cycle?} DFS finds the cycle closed by the first back edge in DFS order; BFS from every start (Approach~1) would find a shortest one. The problem accepts any.
\end{itemize}
\end{tricksbox}

\subsubsection*{5. Trace}

\begin{tracebox}[{Sample: adjacency lists 1: 3,\ 2: 1 4,\ 3: 2 4,\ 4: (none)}]
\begin{center}
\begin{tikzpicture}[x=0.95cm, y=0.95cm, scale=0.92, transform shape]
    \tikzset{gv/.append style={minimum size=6.5mm, font=\scriptsize\bfseries}}
    \foreach \k/\t/\sa/\sb/\sc/\sd in {
        0/{1. enter 1 (gray)}/gvcur/gv/gv/gv,
        1/{2. 1$\to$3 tree; enter 3}/gvseen/gv/gvcur/gv,
        2/{3. 3$\to$2 tree; enter 2}/gvseen/gvcur/gvseen/gv,
        3/{4. 2$\to$1: 1 is gray!}/gvsel/gvcur/gvsel/gv}{
        \begin{scope}[xshift=\k*4.1cm]
            \node[\sa] (t1) at (0,0) {1}; \node[\sc] (t3) at (1.6,0.8) {3};
            \node[\sb] (t2) at (1.6,-0.8) {2}; \node[\sd] (t4) at (3.0,0) {4};
            \draw[gde] (t2) -- (t4); \draw[gde] (t3) -- (t4);
            \ifnum\k>0 \draw[gtree, -{Stealth[length=2.2mm]}] (t1) -- (t3); \else \draw[gde] (t1) -- (t3); \fi
            \ifnum\k>1 \draw[gtree, -{Stealth[length=2.2mm]}] (t3) -- (t2); \else \draw[gde] (t3) -- (t2); \fi
            \ifnum\k>2 \draw[gback] (t2) -- (t1); \else \draw[gde] (t2) -- (t1); \fi
            \node[font=\scriptsize\bfseries] at (1.5,1.55) {\t};
        \end{scope}
    }
\end{tikzpicture}
\end{center}
\begin{center}\small
\begin{tabular}{@{}clll@{}}
\toprule
\textbf{step} & \textbf{at $u$} & \textbf{examine} & \textbf{action} \\
\midrule
1 & 1 & $1 \to 3$, colour 0 & tree, \texttt{par[3] = 1}, stack \texttt{1 3} \\
2 & 3 & $3 \to 2$, colour 0 & tree, \texttt{par[2] = 3}, stack \texttt{1 3 2} \\
3 & 2 & $2 \to 1$, colour 1 & \textbf{back edge}: walk $2, \texttt{par}[2] = 3$, stop at $1$ $\Rightarrow$ \texttt{2 3 1} \\
\bottomrule
\end{tabular}
\end{center}
Reversed: \texttt{1 3 2}; closed: \texttt{1 3 2 1}, $k = 4$. City 4 is never entered. If the adjacency list of 3 had been \texttt{4 2}, DFS would first enter 4, finish it (black), and then the edge $2 \to 4$ would find a \emph{black} vertex --- ignored, correctly.
\end{tracebox}

\subsubsection*{6. Test suite}

\begin{testbox}
\centering\small
\begin{tabular}{@{}p{5.6cm}lp{5.6cm}@{}}
\toprule
\textbf{Input} & \textbf{Our output} & \textbf{Why} \\
\midrule
\texttt{3 3 / 1 2 / 1 3 / 2 3} & \texttt{IMPOSSIBLE} & DAG with a vertex of in-degree 2: the black-target trap \\
\texttt{2 2 / 1 2 / 2 1} & \texttt{3 / 1 2 1} & a 2-cycle is valid here (no parent skip!) \\
\texttt{4 3 / 1 2 / 3 4 / 4 3} & \texttt{3 / 3 4 3} & the cycle is not reachable from 1 \\
\texttt{4 4 / 1 2 / 2 3 / 3 4 / 4 2} & \texttt{4 / 2 3 4 2} & the cycle does not contain the start \\
\texttt{5 5 / 1 2 / 2 3 / 1 3 / 3 4 / 4 5} & \texttt{IMPOSSIBLE} & a forward edge $1 \to 3$ \\
chain $1 \to \dots \to 10^5 \to 1$ & $k = 100001$ & deepest search \\
\bottomrule
\end{tabular}
\end{testbox}
All three approaches were checked on $3000$ random directed graphs with $n \le 9$: a checker validated every printed trip, and ``cycle or not'' was compared with Kahn's algorithm.

\subsubsection*{7. Common mistakes}

\begin{warningbox}[Mistake 1 --- copying the undirected solution]
We ran Problem~\ref{ssec:ch2-round-trip}'s code (``visited and not the parent $\Rightarrow$ cycle'') with one-way adjacency lists. On the first test above it treats the black city 3, reached again by $1 \to 3$, as a cycle; its parent walk from 1 never meets 3, and the program died with \texttt{bad\_alloc}. On the second test it skips the flight $2 \to 1$ as ``the parent edge'' and prints \texttt{IMPOSSIBLE} although $1, 2, 1$ is a valid trip.
\end{warningbox}

\begin{warningbox}[Mistake 2 --- not reversing]
\texttt{2 3 1 2} uses the flights $2 \to 3$ and $3 \to 1$, which do not exist. Reverse the parent walk.
\end{warningbox}

\begin{warningbox}[Mistake 3 --- resetting colours per start]
Re-initialising \texttt{color} inside the loop over $s$ makes every start re-explore the whole graph: \bigO{n(n+m)}, \tle{} on a DAG.
\end{warningbox}

\subsubsection*{8. Takeaways}

\begin{itemize}[leftmargin=*]
    \item Directed cycle $=$ edge to a \emph{gray} vertex. Black vertices are dead ends for cycles.
    \item Undirected and directed cycle detection differ in two details: the parent skip and the meaning of ``visited''. Know both by heart.
    \item An acyclic answer comes with a bonus: the order in which vertices turn black, reversed, is a topological order (Chapter~5).
    \item Kahn's peeling is a DFS-free alternative that both detects and prints a cycle.
\end{itemize}

\begin{variantsbox}
\variant{Only yes/no: can all courses be finished?}{\varnone}
  {report only whether a cycle exists.}
  {the same graph.}
  {either approach without printing; Kahn's version is the usual one (Chapter~5).}
\variant{Contradictory orders}{\varin{https://codeforces.com/problemset/problem/1931/F}}
  {several people report the order of a chat list from their own point of view; is there one order consistent with all reports?}
  {one arrow per pair of consecutive names in a report.}
  {consistent iff the arrow graph has no directed cycle (Exercise E1.5).}
\variant{Which edge to remove}{\varin{https://leetcode.com/problems/redundant-connection-ii/}}
  {a rooted tree plus one extra directed edge; remove an edge so that a rooted tree remains.}
  {either a vertex with two parents, or a directed cycle, or both.}
  {find the in-degree-2 vertex and the cycle; combine the cases (Exercise E1.6).}
\variant{Every city on a cycle}{\varnone}
  {list all cities that lie on some round trip.}
  {a city lies on a cycle iff its strongly connected component has $\ge 2$ cities.}
  {strongly connected components (Chapter~19); for functional graphs the simpler three-state walk of Chapter~1, Pattern~5.}
\end{variantsbox}
